\documentclass[10pt,a4paper]{amsart}
\usepackage{amsmath,amssymb,amsthm,graphicx}
\usepackage{graphicx,color}
\usepackage{euscript}
\usepackage{amsmath}
\usepackage{amsfonts}
\usepackage{amssymb}
\usepackage{amsthm}
\usepackage{hyperref}
\usepackage{epsfig}
\usepackage{epsf}
\usepackage{epstopdf}
\usepackage{caption}
\usepackage{url}
\usepackage{array}
\usepackage{amssymb}\usepackage{amscd}
\usepackage{epic}
\usepackage{eufrak}
\usepackage{tikz,underscore,pgfplots}
\usepackage{tikz-cd}
\usepackage{bm}
\usepackage{float}
\usepackage{anysize}
\usepackage{amsfonts}
\usepackage{pinlabel} 
\usepackage{combelow}
\usepackage{comment}
\newtheorem{theorem}{Theorem}[section] 
\newtheorem{corollary}[theorem]{Corollary} 
\newtheorem{lemma}[theorem]{Lemma} 
\newtheorem{proposition}[theorem]{Proposition} 
\newtheorem{example}[theorem]{Example}
\newtheorem{conjecture}[theorem]{Conjecture}
\newtheorem{problem}[theorem]{Problem}  
\newtheorem{remark}[theorem]{Remark}  
\newtheorem{definition}[theorem]{Definition}
\newtheorem{algorithm}[theorem]{Algorithm}
 

\newcommand{\ch}{\mathrm{char}}

\newcommand{\Kh}{\mathrm{Kh}}

\newcommand{\CA}{\mathcal{A}} 
\newcommand{\CH}{\mathcal{H}}
\newcommand{\Cab}{\mathrm{Cab}}
\newcommand{\col}{\mathrm{col}}
\newcommand{\Tw}{\mathcal{TW}}

\newcommand{\F}{\mathbb{F}}
\newcommand{\E}{\mathbf{E}}
\newcommand{\C}{\mathbb{C}}
\newcommand{\CC}{\mathcal{C}}
\newcommand{\cCFL}{\mathcal{C\!F\!L}}
\newcommand{\cCFK}{\mathcal{C\hspace{-.5mm}F\hspace{-.3mm}K}}
\newcommand{\cHFL}{\mathcal{H\!F\! L}}
\newcommand{\CFL}{\mathit{CFL}}
\newcommand{\HFL}{\mathit{HFL}}
\newcommand{\HFK}{\mathit{HFK}}
\newcommand{\HF}{\mathit{HF}}
\newcommand{\Cone}{\mathrm{Cone}}
\newcommand{\scU}{\mathscr{U}}
\newcommand{\scV}{\mathscr{V}}
\newcommand{\s}{\mathfrak{s}}
\newcommand{\ee}{\mathbf{e}}
\newcommand{\ff}{\mathbf{f}}
\newcommand{\full}{\mathrm{full}}
\newcommand{\sgn}{\mathrm{sgn}}
\newcommand{\lk}{\mathrm{lk}}
\newcommand{\Split}{\mathrm{split}}
\newcommand{\Spin}{\text{Spin}^{c}}
\newcommand{\alphas}{\boldsymbol{\alpha}}
\newcommand{\betas}{\boldsymbol{\beta}}
\newcommand{\gammas}{\boldsymbol{\gamma}}
\newcommand{\deltas}{\boldsymbol{\delta}}
\newcommand{\ovl}{\overline}

\newcommand{\cbHFL}{\boldsymbol{{\mathcal{H\!F\!L}}}}
\newcommand{\cbCFL}{\boldsymbol{\mathcal{C\!F\!L}}}
\newcommand{\bJ}{\boldsymbol{\J}}

\newcommand{\Imm}{\mathrm{Im}}
\newcommand{\cK}{\mathcal{K}}
\newcommand{\bt}{\mathbf{t}}
\newcommand{\HD}{\mathcal{H}}
\newcommand{\G}{\mathcal{G}}
%\newcommand{\i}{\mathit{in}}

\newcommand{\tw}{\mathbf{tw}}
\newcommand{\w}{\mathbf{w}}
\newcommand{\x}{\mathbf{x}}
\newcommand{\y}{\mathbf{y}}

\newcommand{\newword}[1]{\emph{\textbf{#1}}}

\newcommand{\spinc}{\text{Spin}^{c}}


\newcommand{\Alt}{\mathrm{Alt}}
\newcommand{\HY}{\mathrm{HY}}

\def\L{\mathcal{L}}
\def\gr{\textup{gr}}
\def\w{{\bf{w}}}
\def\z{{\bf{z}}}
\def\d{\mathbf{d}}
\def\x{\mathbf{x}}
\def\y{\mathbf{y}}

\def\k{\mathbf{k}}
\def\m{\mathbf{m}}
\def\bF{\mathcal{F}}


\newcommand{\HHH}{\mathrm{HHH}}

\newcommand{\FT}{\mathrm{FT}}

\newcommand{\HH}{\mathbb{H}}
\newcommand{\Z}{\mathbb{Z}}
\newcommand{\R}{\mathbb{R}}
\newcommand{\UU}{\mathbf{U}}
\newcommand{\A}{\mathfrak{A}}
\newcommand{\CT}{\mathcal{T}}

\newcommand{\Br}{\mathrm{Br}}

\newcommand{\J}{\mathcal{J}}

\renewcommand{\AA}[1]{{\color{blue}#1}}
\newcommand{\EG}[1]{{\color{red}#1}}
\newcommand{\BL}[1]{{\color{purple}#1--Beibei}}
%\def\aa#1{{\textcolor{blue}{#1}}}

\newcommand{\spl}{\mathrm{sp}}

\newcommand{\cc}{\mathbf{c}}

\newcommand{\ttt}{\mathbf{t}}

\newcommand{\GCD}{\mathrm{GCD}}
\newcommand{\CFK}{\mathrm{CFK}}
\newcommand{\CF}{\mathcal{F}}

\newcommand{\Id}{\mathrm{Id}}

\newcommand{\spann}{\mathrm{span}}

\newcommand{\stab}{\mathrm{stab}}

\renewcommand{\ss}{\mathbf{s}}

\newcommand{\AAA}{\mathsf{A}}
\title{Cable algebra action and rank-one freeness of unknot cables}
\author{Akram Alishahi, Eugene Gorsky, Beibei Liu}
\date{arxiv:2508.21776v1; CC BY 4.0}
\begin{document}
\maketitle
\noindent\textit{Editorial scope.} Three source index slips in the supporting torus-link calculation are preserved verbatim. In original source line 1874 the product index is inconsistent; compare the explicit multi-index convention at 1938. Line 1941 repeats the h-value with qm+r; the author has already derived qm+2r at 1909 and 1931 and the increment q+1 at 1937. At 1985 the printed V\_(b-1) denotes the module generator written Y\_(b-1) in the preceding equality. These readings point to formulas already present in the author proof; no mathematical formula in this package has been edited.\par\medskip
\noindent\textit{Editorial note.} Three source index slips in the supporting torus-link calculation are preserved verbatim. In original source line 1874 the product index is inconsistent; compare the explicit multi-index convention at 1938. Line 1941 repeats the h-value with qm+r; the author has already derived qm+2r at 1909 and 1931 and the increment q+1 at 1937. At 1985 the printed V\_(b-1) denotes the module generator written Y\_(b-1) in the preceding equality. These readings point to formulas already present in the author proof; no mathematical formula in this package has been edited.\par\medskip
\noindent\textit{Allowed standard prerequisites (editorial).}\begin{itemize}
\item BLZ Lemma7.1 and GH Theorem4.3. L-space h-function presentation and cable h formula; retained actual citation and bibliography
\item GH and AGL; prop:gradingshifts851–860. L-space property and injectivity of negative-definite cobordisms, and precise grading shifts
\end{itemize}\medskip
\subsection{Heegaard Floer homology}

In this subsection, we review the necessary background for link Floer homology. We assume some familiarity with the basic Heegaard Floer homology and its refinement to knots, see \cite{OSProperties, OSKnots, OSLinks, RasmussenKnots} for details.  Given a link $\L$ with $n$ components in a 3-manifold $Y$, one can define the ``full" version of Heegaard Floer complex $\cCFL(Y, L)$ from the Heegaard diagram by using some version of Lagrangian Floer complexes. The multi-pointed Heegaard link diagram $(\Sigma, \alphas, \betas,\w, \z)$ for such pair $(Y, L)$ is defined as follows:
\begin{enumerate}
    \item  $\Sigma$ is a closed oriented genus $g$ surface;
    \item $\alphas=\{ \alpha_1, \cdots, \alpha_{g+n-1} \}$, and $\betas=\{ \beta_1, \cdots, \beta_{g+n-1}\}$ are collections of simple closed curves on $\Sigma$ such that $\alphas$ and $\betas$ each span a $g$-dimensional subspace of $H_1(\Sigma, \Z)$ and  curves in $\alphas$ (and $\betas$) are pairwise disjoint.
    \item $\w=\{w_1, \cdots, w_n\}$ and $\z=\{z_1, \cdots, z_n\}$ denote sequences of basepoints on $\Sigma$ such that each component of $\Sigma\setminus \alphas$ (respectively of $\Sigma\setminus \betas$) contains a single point of $\w$ and a single point of $\z$. 
\end{enumerate}

%Any two Heegaard link diagram for a pair $(Y, L)$ can be connected by a sequence of Heegaard moves for link diagrams. See \cite[Theorem 4.7]{OSLinks}. 
The generators of the link Floer chain complex $\cCFL(Y, L)$ are tuples of intersection points $\x\in \mathbb{T}_\alpha\cap \mathbb{T}_{\beta}$ of the Lagrangian tori
$$\mathbb{T}_\alpha=\alpha_1\times \cdots \times \alpha_{g+n-1}, \quad \mathbb{T}_{\beta}=\beta_1\times \cdots \times \beta_{g+n-1}$$
in the symmetric product $\textup{Sym}^{g+n-1}(\Sigma)$. 
 
We work over the coefficient field $\F=\Z/ 2\Z$ in the paper. The marked points $z_i$ and $w_i$ are associated with formal variables $V_i$ and $U_i$, and the chain complex $\cCFL(Y,L)$ is a module over  the polynomial ring $R=\F[U_1, \cdots, U_n, V_1, \cdots, V_n]$. The differential in $\cCFL(Y, L)$ is obtained by counting pseudo-holomorphic disks in $\textup{Sym}^{g+n-1}(\Sigma)$ via:
\begin{equation}
\label{eq: HF differential}
\partial \x=\sum_{\y\in \mathbb{T}_\alpha\cap \mathbb{T}_\beta}\sum_{\varphi\in \pi_2(\x, \y),\mu(\varphi)=1}(\# \mathcal{M}(\varphi)/\R) U_1^{n_{w_1}(\varphi)}\cdots U_n^{n_{w_{n}}(\varphi)}V_1^{n_{z_{1}}(\varphi)}\cdots V_n^{n_{z_n}(\varphi)} \y. 
\end{equation}
Here the sum is taken over all holomorphic disks from $\x$ to $\y$ in $\pi_2(\x, \y)$ representing the homotopy classes of maps $\varphi$ from the unit disk $\mathbb{D}\subset \mathbb{C} $ to $\textup{Sym}^{g+n-1}(\Sigma)$ satisfying some boundary condition.  The Maslov index of $\varphi$ is denoted by $\mu(\varphi)$. For any point $x\in \Sigma\setminus (\alphas\cup \betas)$, $n_x(\varphi)$ denotes the intersection number of the divisor $\{x\} \times \textup{Sym}^{g+n-2}(\Sigma)\subset \textup{Sym}^{g+n-1}(\Sigma)$ with $\varphi(\mathbb{D})$.  The moduli space $\mathcal{M}$ consists of all pseudo-holomorphic curves representing $\varphi$ for a generic 1-parameter family of almost complex structures on $\textup{Sym}^{g+n-1}(\Sigma)$. For more details, see \cite{OSProperties}. 

We denote the homology of the link Floer chain complex $\cCFL(Y, L)$ by $\cHFL(Y, L)$, which is an invariant of the pair $(Y, L)$.  The actions of $U_iV_i$ on the complex $\cCFL(L)$ are pairwise homotopic, so we introduce the following relation:
$$U_1V_1=\cdots =U_nV_n=\UU,$$
and let $R_{UV}$ denote the ring generated by $U_1, \cdots, U_n, V_1, \cdots, V_n$ modulo this relation.  Then the action of $R$ on $\cHFL(Y, L)$ factors through $R_{UV}$. We get different versions of Heegaard Floer homology by setting extra constraints on the  variables $U_i, V_i$. For example, $H_\ast(\cCFL(Y, L)/(V_1-1, \cdots, V_n-1))$ is isomorphic to $\mathit{HF}^-(Y)$, the minus version of Heegaard Floer homology of the 3-manifold $Y$. 

In this paper, we focus on links $L=L_1\cup \ldots \cup L_n$ in the three-sphere and, for simplicity, we write the complex as $\cCFL(L)$. The complex $\cCFL(L)$ and its homology $\cHFL(L)$ have several gradings. Firstly, there is the $\mathbb{Q}\times \mathbb{Q}$-valued Maslov bigrading, denoted by $(\gr_{\w}, \gr_{\z})$, as well as the Alexander multi-grading $A=(A_1, \ldots, A_n)$ valued in the lattice
$$\HH_L=\Z^n+\dfrac{1}{2}(\ell_1, \ldots, \ell_n)$$
where $\ell_i$ is the linking number of $L_i$ with the rest of the components, i.e., $\ell_i=\sum_{j\neq i} lk(L_i, L_j)$. These gradings are relatively determined by the following equations:
\begin{equation}\label{eq:relgradings}
\begin{split}
&\gr_{\w}(\x)-\gr_{\w}(\y)=\mu(\varphi)-2n_{\w}(\varphi)\\
&\gr_{\z}(\x)-\gr_{\z}(\y)=\mu(\varphi)-2n_{\z}(\varphi)\\
&A_i(\x)-A_{i}(\y)=n_{z_i}(\varphi)-n_{w_i}(\varphi)
\end{split}
\end{equation}
where $\varphi\in\pi_2(\x,\y)$ and $n_{\w}(\varphi)=\sum_{i=1}^nn_{w_i}(\varphi)$ whereas $n_{\z}(\varphi)=\sum_{i=1}^nn_{z_i}(\varphi)$. Furthermore, the actions of $U_i,V_i$ are homogeneous with respect to these gradings with weights
\begin{equation}
\label{eq: U V gradings}    
(\gr_{\w}, \gr_{\z})(U_i)=(-2, 0), \quad (\gr_{\w}, \gr_{\z})(V_i)=(0, -2), \quad \text{and } A(V_i)=-A(U_i)=\ee_i
\end{equation}
where $\ee_i$ is the standard $i$-th coordinate vector in $\R^n$. 
The differential on $\cCFL(L)$ preserves the Alexander grading, while dropping both $\gr_{\w}$ and $\gr_{\z}$ by  one. In fact, these gradings satisfy the relation
$$A_1+\ldots+A_n=\frac{1}{2} (\gr_{\w}-\gr_{\z}).$$



%The differential on $\cCFL(L)$ preserves the Alexander grading, which is equivalent to the following:\EG{Reference, does not need a proposition; also add equations for $\gr_{\w}$ and $\gr_{\z}$}
%\begin{proposition} 
%\label{prop: Alexander from disks}
%Suppose $\varphi\in \pi_2(\x,\y)$ is a holomorphic disk from $\x$ to $\y$. Then 
%$$
%A_i(\x)-A_i(\y)=n_{z_i}(\varphi)-n_{w_i}(\varphi).
%$$
%\end{proposition}



For any $\ss\in \HH_L$, let $\cCFL(L, \ss)$ denote the subcomplex of $\cCFL(L)$ generated by all elements $x$ such that $A(x)=\ss$. By the grading formula \eqref{eq: U V gradings}, it is straightforward to see that the product $U_iV_i$ preserves the Alexander gradings, and $\cCFL(L, \ss)$ is a module over the subring $\F[U_1V_1, U_2V_2, \cdots, U_nV_n]$, and so $\cHFL(L, \ss)$ is an $\F[\UU]$-module. In particular, for any link $L$ we have
$$\cHFL(L)=\oplus_{s\in \HH_L} \cHFL(L, \ss).$$
and $\cHFL(L, \ss)$ is the direct sum of one copy of $\F[\UU]$ with some $\UU$-torsion. 



Recall that a rational homology sphere $M$ is an \emph{L-space} if for each $\Spin$-structure $\mathfrak{s}$,  $HF^-(M, \mathfrak{s})$ is isomorphic to $\F[\UU]$. A link $L$ in the three-sphere $S^3$ is an \emph{L-space link} if the Dehn surgery $S^3_{\mathbf{d}}(L)$ is an L-space for all $\mathbf{d}\gg 0$. Another way to characterize the L-space links is the condition that $\cHFL(L, \ss)\cong \F[\UU]$ for every lattice point $\ss\in \HH_L$. For any link $L$ in the three sphere, we define the link invariant,  known as the $h$-function.

\begin{definition}
    The $h$-function $h:\HH_L\to \Z$ of a link $L$ in the three-sphere is defined so that  for a given $\ss\in \HH_L$ the value $h(\ss)$ is  $-\dfrac{1}{2} \gr_{\w}$ for the maximal Maslov grading of non-torsion elements in $\cHFL(L, \ss)$. 
\end{definition}

For L-space links $L$, its $h$-function can be computed from the Alexander polynomials of the link and its sublinks, see \cite{BorodzikGorskyImmersed}. Moreover, for such links, $\cHFL(L)$ is determined by its $h$-function (see \cite{BLZ}), which is determined by these Alexander polynomials. 

Let $L$ be an $n$-component link in the three-sphere, and $\overline{L}$ be another link obtained from $L$ by 
adding a full twist on $n$ strands, that is, $\overline{L}$ is  obtained from $L$ by blowing down an $(-1)$-framed unknot circling the $n$ strands of $L$ where we assume that all strands are oriented the same way. Consider the corresponding cobordism from $(S^3, L)$ to $(S^3, \overline{L})$ obtained by attaching a $2$-handle to $S^3\times I$ along the aforementioned $(-1)-$framed unknot in $S^3\times\{1\}$. Let $[S^2]$ denote the generator of the second homology class of the cobordism, and let $\phi_k:\cHFL(L)\to \cHFL(\overline{L})$ denote the induced cobordism maps in homology in the $\spinc$-structure $\mathfrak{s}_k$ such that  $\langle \mathfrak{s}_k, [S^2]\rangle =2k+1$.  See \cite{AGL} for a detailed discussion of the properties of $\phi_k$.  We have

\begin{proposition}\cite{AGL}\label{prop:gradingshifts}
Given an n-component link $L$ in the three sphere, the cobordism maps $\phi_k:\cHFL(L)\to\cHFL(\overline{L}) $ induced from the $(-1)$-surgery on the unknot with $0\leq k\leq n-1$ satisfy the following grading properties:
$$
\gr_{\w}(\phi_k)=-k^2-k,\ A_i(\phi_k)=-k+(n-1)/2,
$$
where $i=1, \cdots, n$. 
\end{proposition}

\section{Colored knot Floer homology of the unknot}

In this section, we compute the colored knot Floer homology of the unknot.  %From now on we will fix $n$ and denote the cobordism maps $\phi_k$ from Proposition \ref{prop:gradingshifts} simply by $\phi_k$.


\subsection{Homology of $T(n,mn)$} The first step is computing the link Floer homology of the torus link $T(n,nm)$. Since $m$ is fixed in this subsection, we will denote $\cc_m=\frac{m(n-1)}{2}$ by $\cc$ to simplify notations. 

\begin{theorem}
\label{thm: Tnm}
The homology $\cHFL(T(n,mn))$ as a module over $R_{UV}$ is generated by $Y_0,\ldots,Y_{m(n-1)}$ with Alexander and Maslov degrees 
$$
A(Y_i)=(\cc-i,\ldots,\cc-i),\ \gr_{\w}(Y_i)=-(q+1)(qm+2r),\ \ i=qm+r,\ 0\le r\le m-1.
$$
Furthermore, the relations between $Y_i$ are spanned by $U_IY_i=V_{\overline{I}}Y_{i+1}$ for all subsets $I\subset\{1,2,\cdots,n\}$ with $q+1$ elements. Here, $\overline{I}=\{1,2,\cdots,n\}\setminus I$, $U_I=\prod_{j\in I}U_i$ and $V_{\overline{I}}=\prod_{j\in \overline{I}}V_j$.
\end{theorem}





%\begin{theorem}
%\label{thm: Tnm}
%The homology $\cHFL(T(n,mn))$ has $m(n-1)+1$ generators $Y_0,\ldots,Y_{m(n-1)}$ of degrees 
%$$
%A(Y_i)=(\cc-i,\ldots,\cc-i),\ \gr_{\w}(Y_i)=-(q+1)(qm+2r),\ \ i=qm+r,\ 0\le r\le m-1.
%$$
%Here $\cc=\frac{m(n-1)}{2}$.
%\end{theorem}

\begin{proof}
The proof is similar to the proof of \cite[Theorem 7.3]{BLZ}. By  \cite[Lemma 7.1]{BLZ}, to write down the generators we need to compute the $h$-function for $T(n,mn)$.  Since the Alexander polynomial for $T(n,mn)$ is symmetric in $t_1,\ldots,t_n$, the $h$-function is symmetric in $s_1,s_2,\cdots,s_n$ so we consider the case $s_1\le s_2\le\cdots \le s_n$. By \cite[Theorem 4.3]{GH} for $s_1\le s_2\le\cdots \le s_n$ we get
\begin{equation}
\label{eq: h function cable}
h(s_1,\ldots,s_n)=h_0(s_1-\cc)+h_0(s_2-\cc+m)+\ldots+h_0(s_n-\cc+m(n-1))
\end{equation}
where  
$$
h_0(x)=\begin{cases}
0 & \mathrm{if}\ x\ge 0\\
-x & \mathrm{if}\ x<0
\end{cases}
$$
In particular, for $\cc\le s_1\le \ldots\le s_n$ we get $h(s_1,\ldots,s_n)=0$. Suppose that $s=\cc-i$ where $i=mq+r\ge 0$ then (assuming $i\le m(n-1)=2\cc)$ we get
$$
s-\cc\le 0, s-\cc+m\le 0,\ldots,s-\cc+qm\le 0\ \mathrm{but}\ s-\cc+(q+1)m>0
$$
so 
$$
h(s,\ldots,s)=(q+1)i-\frac{1}{2}q(q+1)m=\frac{1}{2}(q+1)(qm+2r).
$$
We claim that 
\begin{equation}
\label{eq: change h diagonal}
h(s-1,\ldots,s-1)-h(s,\ldots,s)=q+1.
\end{equation}
 If $r<m-1$ then going from $s$ to $s-1$ changes $i$ to $i+1$ and $r$ to $r+1$ while preserving $q$, so \eqref{eq: change h diagonal} is clear. If $r=m-1$ then $i=qm+(m-1),i+1=(q+1)m+0$ and
$$
h(s,\ldots,s)=\frac{1}{2}(q+1)(qm+2m-2),\ h(s-1,\ldots,s-1)=\frac{1}{2}(q+2)(q+1)m
$$
hence \eqref{eq: change h diagonal} holds as well.

Furthermore, one can check that
$$
h(\underbrace{s-1,\ldots,s-1}_{k},s,\ldots,s)=
\begin{cases}
h(s,\ldots,s)+k & k\le q+1\\
h(s-1,\ldots,s-1)=h(s,\ldots,s)+q+1 & k\ge q+1.
\end{cases}
$$
Let $Y_i$ be the generator of the $\F[\UU]$-tower in Alexander grading $(\cc-i,\ldots,\cc-i)$. It will have Maslov grading  
\[-2h(\cc-i,\ldots,\cc-i)=-(q+1)(qm+2r)\]
and by \cite[Lemma 7.1]{BLZ}, $Y_0,Y_1,\cdots,Y_{m(n-1)}$ generate $\cHFL(T(n,mn))$.

By \cite[Lemma 7.1]{BLZ}, the relations are spanned by the relations between $Y_a$ and $Y_b$ for any $0\leq a\leq b \leq (n-1)m$ which are the ones spanned by $(U_iV_i+U_jV_j)Y_a$ and $(U_iV_i+U_jV_j)Y_b$ as well as sums
$$U^I Y_a+V^J Y_b$$
ranging over sequences of nonnegative integer $I=(I_1,\cdots,I_n)$ and $J=(J_1,\cdots,J_n)$ such that 
$$I+J=(b-a, \cdots, b-a) \text{ and } |I|_{L^1}:=\sum_{k=1}^n I_k=(q_b+1)(q_b m+2r_b)/2-(q_a+1)(q_a m+2r_a)/2$$
where $a=q_a m+r_a$ and $b=q_b m+r_b$. Here, $U^I=\prod_{k}U_k^{I_k}$ and $V^J=\prod_kV_k^{J_k}$.


If $a=i$ and $b=i+1$, then $(q_b+1)(q_b m+r_b)/2-(q_a+1)(q_a m+r_a)/2=q+1$ and so these relations are exactly the ones described in the statement.



% The relation between $Y_k$ and $Y_{k+1}$ can be observed from the computation of its Alexander polynomial and its $h$-function that
% $$h(\cc-k-1, \cdots, \cc-k-1)-h(\cc-k, \cdots, \cc-k)=q+1.$$
% By \cite[Lemma 7.1]{BLZ}, the relations in $\cHFL(T(n, mn))$ is spanned by the relations between $Y_a$ and $Y_b$ for any $0\leq a\leq b \leq (n-1)m$ which are the ones spanned by $(U_iV_i+U_jV_j)Y_a$ and $(U_iV_i+U_jV_j)Y_b$ as well as sums
% $$U^I Y_a+V^J Y_b$$
% ranging over sequences of nonnegative integer $I$ and $J$ such that 
% $$I+J=(b-a, \cdots, b-a) \text{ and } |I|_{L^1}=(q_b+1)(q_b m+r_b)/2-(q_a+1)(q_a m+r_a)/2$$
% where $a=q_a m+r_a$ and $b=q_b m+r_b$. 

Suppose $b>a+1$. We show that the relations between $Y_a$ and $Y_b$ are in the span of the relations between consecutive $Y_i$ and $Y_{i+1}$. We prove this claim by induction on $b$. First, we show that there is a tuple $I'=(I'_1,\cdots,I'_{n})$ such that for all $1\leq i\leq n$, $0\leq I'_i\leq I_i$, and
$$|I'|_{L^1}=q_{b-1}+1, \quad |I'|_{L^\infty}=\max\{I'_1,\cdots,I'_{n}\}\leq 2 \text{ and } |I-I'|_{L_{\infty}}\le b-a-1$$

For $k\in \mathbb{N}$, we write
$$a_k:=\# \{i: I_i\geq k\}.$$
Then 
$$\sum_{k=1}^{b-a} a_k=|I|_{L^1}=(q_b+1)(q_b m+2r_b)/2-(q_a+1)(q_a m+2r_a)/2, \quad \text{ and } a_1\geq a_2\geq \cdots \geq a_{b-a},$$
and so by Equation (\ref{eq: change h diagonal}) $\sum_{k=1}^{b-a}a_k=\sum_{i=a}^{b-1}(q_i+1)$. Consequently, the number of entries $I_i$ that are equal to $b-a$ is at most $q_{b-1}+1$. So, such $I'$ exists if and only if $a_1+a_2\geq q_{b-1}+1$. Suppose to the contrary that 
$a_1+a_2< q_{b-1}+1$. Then $a_2<(q_{b-1}+1)/2$ and therefore

$$\sum_{i=a}^{b-1}(q_i+1)=\sum_{k=1}^{b-a} a_k=a_1+a_2+\sum_{k=3}^{b-a}a_k< q_{b-1}+1+(b-a-2) \dfrac{q_{b-1}+1}{2}=(b-a)\frac{q_{b-1}+1}{2}.$$

On the other hand, $q_{b-1}+1\le q_{b+a-1}+1\le (q_{a+i}+1)+(q_{b-i-1}+1)$ implies that 
\[\sum_{i=a}^{b-1}(q_i+1)\ge \frac{b-a}{2}(q_{b-1}+1)\]
which is a contradiction. 









Therefore such a tuple $I'$ exists. By induction, for $J'=(b-a-1, \cdots, b-a-1)-(I-I')$, the relation 
$$U^{I-I'}Y_a=V^{J'} Y_{b-1}$$
is in the span of the claimed relations between consecutive generators. Then, this relation implies that 
$$U^I Y_a=U^{I'} U^{I-I'}Y_a=U^{I'} V^{J'}Y_{b-1}.$$

If $|I'|_{L^\infty}=1$, the relations between $Y_{b-1}$ and $Y_b$ imply that
$$U^I Y_a=U^{I'} V^{J'}Y_{b-1}=V^J Y_{b}.$$
Otherwise, for any  index $i$ with  $I'_i=2$, we observe that $(I-I')_i<b-a-1$, hence $J'_i>0$. Then $U^{I'} V^{J'}$ contains a factor of $U_iV_i$. Since, $|I'_i|_{L_1}=q_{b-1}+1\le n$ there is an index $j$ with $I'_j=0$. We will trade the $U_iV_i$ factor with $U_jV_j$.  By repeating this process enough times, we obtain another tuple $I''$ from $I'$ such that $|I''|_{L^1}=|I'|_{L^1}$ and $|I''|_{L^\infty}=1$. That is
$$U^{I'} V^{J'} Y_{b-1}=U^{I''}V^{J''}V_{b-1}=V^{J} Y_{b}.$$
\end{proof}

\subsection{Cable algebra and cobordism maps} Next, we explicitly describe the cobordism maps $\phi_k$ corresponding to adding full twists  for $T(n,mn)$. In order to do that, we consider the direct sum of homology for all $n$-strand cables of the unknot with nonnegative integer slopes:
\begin{equation}\label{eq:CabU}
\Cab_n(O): =\Tw(O_n)=\bigoplus_{m=0}^{\infty}\cHFL(T(n,nm))
\end{equation}
where $O$ is the unknot and  $O_n$ is the $n$-component unlink. 

Each summand is $\Z^n\oplus \Z$ graded by Alexander multi-grading ($\Z^n$) and Maslov grading ($\Z$). Moreover, $\Cab_n(O)$ has an additional $\Z$ grading given by the cabling slope $m$, we call it \emph{twist grading} and denote it by $\tw$. Note that cobordism maps $\phi_k$, corresponding to adding a full twist, act naturally on $\Cab_n(O)$ with twist grading $1$ i.e. they map $\cHFL(T(n,mn))$ to $\cHFL(T(n,(m+1)n))$. 
%We would like to explicitly describe the action of the cobordism maps $\phi_k$ on $\Cab_n(O)$. 

\begin{definition}
The $n$-strand \newword{cable algebra} $\CA_n$ is defined as a $\Z^n\oplus \Z\oplus \Z$ graded algebra  over $R_{UV}$ with commuting generators $a_0,\ldots,a_{n-1}$ and the following relations:
\begin{itemize}
\item Linear relations
\begin{equation}
\label{eq: linear}
U_Ia_{k-1}=V_{\overline{I}}a_k
\end{equation}
where $I\subset \{1,\ldots,n\}$ is an arbitrary subset with $|I|=k$ and $\overline{I}=\{1,\ldots,n\}\setminus I.$ 
%\AA{Change $K$ to $I$, because $K$ i the knot}\EG{done}
\item Quadratic relations
\begin{equation}
\label{eq: quadratic}
a_ia_j=\UU^{k\ell-ij}a_{k}a_{\ell}
\end{equation} 
whenever $i+j=k+\ell$ and $i\le k\le \ell\le j$.
\end{itemize}
As above, the gradings correspond to the Alexander ($\Z^n$) and Maslov ($\Z$) gradings, along with an additional  $\Z$ \emph{twist grading}. The generator $a_k$ has Alexander grading 
$$
A(a_k)=\left(\frac{n-1}{2}-k,\ldots, \frac{n-1}{2}-k\right),
$$
Maslov grading $\gr_{\w}(a_k)=-k^2-k$ and twist grading $\tw(a_k)=1$.
\end{definition}

It is easy to check that the relations are homogeneous with respect to all three gradings.

\begin{theorem}
\label{thm: cabled algebra}
The homology $\Cab_n(O)$ is a $\Z^n\oplus \Z\oplus \Z$ graded module over the algebra $\CA_n$ where the generators $a_k$ act by cobordism maps $\phi_k$. Furthermore, $\Cab_n(O)$ is a free $\CA_n$-module generated by a single class $1\in \cHFL(T(n,0))$.
\end{theorem}

\begin{proof}
We proceed in several steps.

{\bf Step 1:} We show that there is a well-defined action of algebra $\CA_n$ on $\Cab_n(O)$, where $a_k$ acts by the cobordism map $\phi_k$. First, we prove that the actions of $\phi_{k}$ and $\phi_{k'}$ on $\Cab_n(O)$ commute. By \cite{GH} all links $T(n,nm)$ are L-space links, so $\cHFL(T(n,mn))$ is free of rank 1 over $\F[\UU]$ in each Alexander degree. The maps $\phi_{k}$ and $\phi_{k'}$ correspond to negative definite cobordisms, so they are both injective \cite{AGL}, and hence determined by their Alexander and Maslov degrees. Now we observe that both compositions $\phi_{k}\circ \phi_{k'}$ and $\phi_{k'}\circ \phi_{k}$
have the same Alexander and Maslov degrees by Proposition \ref{prop:gradingshifts}, and hence coincide.

Next, we check that these cobordism maps satisfy  relations \eqref{eq: linear} and \eqref{eq: quadratic}. Note that $a_k$ has the same Alexander, Maslov and twist grading as $\phi_k$, thus relations \eqref{eq: linear} and \eqref{eq: quadratic} being homogeneous, implies that the cobordism maps satisfy  these relations as well and the action is well-defined. 

% This allows us to define the action of commuting operators $a_k$ on $\Cab_n(O)$. We need to prove that they satisfy the relations \eqref{eq: linear} and \eqref{eq: quadratic}. This can be done again by comparing the Alexander and Maslov degrees for both sides. For \eqref{eq: linear}, this follows from \cite[Theorem 7.3]{BLZ}. \AA{Is this the correct citation here?}

% For \eqref{eq: quadratic}, observe that the monomials $a_ia_j$ and $a_{k}a_{\ell}$ have the same Alexander degree whenever $i+j=k+\ell$. On the other hand, their Maslov gradings differ by
% $$
% i(i-1)+j(j-1)-k(k-1)-\ell(\ell-1)=i^2+j^2-k^2-\ell^2-(i+j)+(k+\ell)=
% $$
% $$
% (i+j)^2-(k+\ell)^2-(i+j)+(k+\ell)+2(k\ell-ij)=2(k\ell-ij).
% $$
% Note that the inequality $i\le k\le \ell\le j$ ensures $k\ell-ij\ge 0$.

% To sum up, we proved that the algebra $\CA_n$ acts on $\Cab_n(O)$ via the cobordism maps $\phi_k$.

{\bf Step 2:} We prove that $\Cab_n(O)$ is generated by a single class $1\in \cHFL(T(n,0))$ under the action of $\CA_n$. 
We follow Theorem \ref{thm: Tnm} and claim that the generator $Y_i\in \cHFL(T(n,mn))$ can be identified with the monomial $a_{q}^{m-r}a_{q+1}^{r}$ (where $i=mq+r$) applied to $1$.
By the above, the element $a_{q}^{m-r}a_{q+1}^{r}(1)$ is nonzero, and it is sufficient to verify that it has the same Alexander and Maslov degree as $Y_i$. Indeed, the Alexander grading equals
$$
A_j(a_{q}^{m-r}a_{q+1}^{r})=(m-r)\left(\frac{n-1}{2}-q\right)+r\left(\frac{n-1}{2}-q-1\right)=m\frac{n-1}{2}-mq-r=\cc-i
$$
and the Maslov grading equals
$$
\gr_{\w}(a_{q}^{m-r}a_{q+1}^{r})=-(m-r)q(q+1)-r(q+1)(q+2)=-(q+1)(mq+2r).
$$
To sum up, we have a surjective map of $\CA_n$ modules from $\CA_n$ to $\Cab_n(O)$.

{\bf Step 3:} Finally, we need to prove that this map is an isomorphism and $\Cab_n(O)\simeq \CA_n$. For this, we need to find the generators of $\CA_n$ over $R_{UV}$, relations between them and compare these with Theorem \ref{thm: Tnm}. 

Define $\widetilde{Y_i}=a_{q}^{m-r}a_{q+1}^{r}\in \CA_n$ where $i=mq+r$. Let us prove that $\widetilde{Y_i}$ generate $\CA_n$ over $R_{UV}$. Indeed, consider a monomial $M=a_{i_1}\cdots a_{i_m}$ with $i_1\le i_2\le \cdots \le i_m$. If $i_m-i_1\le 1$ then $M=\widetilde{Y}_{i_1+\ldots+i_m}$ by the above. Otherwise we can use \eqref{eq: quadratic} to write
$$
a_{i_1}a_{i_m}=\UU^{i_m-i_1-1}a_{i_1+1}a_{i_m-1}.
$$ 
since $(i_1+1)(i_m-1)-i_1i_m=i_m-i_1-1$.
This decreases the power of $a_{i_1}$ and we can proceed by induction.

Finally, we need to verify that the relations from Theorem \ref{thm: Tnm}  hold for $\widetilde{Y_i}$. This is clear since the relations between
$\widetilde{Y}_i=a_{q}^{m-r}a_{q+1}^{r}$ and $\widetilde{Y}_{i+1}=a_{q}^{m-r-1}a_{q+1}^{r+1}$ agree with the relations between $a_{q}$ and $a_{q+1}$,
up to multiplication by $a_{q}^{m-r-1}a_{q+1}^{r}$. 
\end{proof}

\begin{thebibliography}{99}

\bibitem{AGL} A. Alishahi, E. Gorsky, B. Liu. Splitting maps in link Floer homology and integer points in permutahedra. arXiv:2307.07741

\bibitem{BGHW} W. Ballinger, E. Gorsky, M. Hogancamp, J. Wang. Stable deformed $\mathfrak{gl}_N$ homology of torus knots. arXiv:2507.00175


\bibitem{BPRW}  A. Beliakova, K. Putyra, L.-H. Robert, E. Wagner. A proof of Dunfield-Gukov-Rasmussen Conjecture. J. Eur. Math. Soc. (2025), arXiv:2210.00878

\bibitem{BorodzikGorskyImmersed} M. Borodzik and E. Gorsky. Immersed concordances of links and Heegaard Floer homology. Indiana Univ. Math. J. 67 (2018), no. 3, 1039--1083. 



\bibitem{BLZ} M. Borodzik, B. Liu, I. Zemke. Lattice homology, formality,  and plumbed L-space links.  J. Eur. Math. Soc. (2024), arXiv:2210.15792

\bibitem{Chen} D. Chen, I. Zemke, H. Zhou. L-space satellite operators and knot Floer homology. arXiv:2412.05755

\bibitem{Cautis} S. Cautis. Clasp technology to knot homology via the affine Grassmannian. Mathematische Annalen 363 (2015): 1053--1115.

\bibitem{Conners} L. Conners. Row--column mirror symmetry for colored torus knot homology.
Selecta Mathematica 30 (5), article 97.

\bibitem{Cooper} B. Cooper, R. Deyeso.  Holonomicity from a Heegaard-Floer Perspective. arXiv:2501.01519

\bibitem{CK} B. Cooper, V. Krushkal. Categorification of the Jones--Wenzl projectors. Quantum Topology 3.2 (2012): 139--180.

\bibitem{DGR} N. Dunfield, S. Gukov, J. Rasmussen. The superpolynomial for knot homologies. Experimental Mathematics 15.2 (2006): 129--159.

\bibitem{GHog} E. Gorsky, M. Hogancamp. Hilbert schemes and  $y$-ification of Khovanov-Rozansky homology.
Geom. Topol. 26 (2022), no. 2, 587--678.

\bibitem{GH} E. Gorsky, J. Hom. Cable links and L-space surgeries.  
Quantum Topology 8 (2017), no.4, 629--666.

\bibitem{GNR} E. Gorsky, A. Negu\cb{t}, J. Rasmussen. Flag Hilbert schemes, colored projectors and Khovanov-Rozansky homology. Advances in Mathematics 378 (2021): 107542.

\bibitem{GOR} E. Gorsky, A. Oblomkov, J. Rasmussen. On stable Khovanov homology of torus knots. Experimental Mathematics 22.3 (2013): 265--281.

\bibitem{Hedden} M. Hedden. On knot Floer homology and cabling
Algebraic \& Geometric Topology 5 (2005), 1197--1222.

\bibitem{Hedden2}  M. Hedden. On knot Floer homology and cabling. II.
International Mathematics Research Notices, 2009, no. 12, 2248--2274.

\bibitem{Hogancamp} 
M. Hogancamp. Categorified Young symmetrizers and stable homology of torus links.
Geometry \& Topology 22 (2018) 2943--3002.

\bibitem{Hogancamp2} M. Hogancamp. A polynomial action on colored  $\mathfrak{sl}_2$  link homology.
Quantum Topol. 10 (2019), no. 1, 1--75.


\bibitem{Kh}  M. Khovanov. A categorification of the Jones polynomial.
Duke Math. J. 101 (2000), no. 3, 359--426.

\bibitem{KR1} M. Khovanov, L. Rozansky. Matrix factorizations and link homology.
Fund. Math. 199 (2008), no. 1, 1--91.

\bibitem{KR2}  M. Khovanov, L. Rozansky. Matrix factorizations and link homology. II
Geom. Topol. 12 (2008), no. 3, 1387--1425.


\bibitem{Licata} J. Licata. Heegaard Floer homology of $(n,n)$-torus links: computations and questions.  arXiv:1208.0394.

\bibitem{Lipshitz-cylindrical} R. Lipshitz. A cylindrical reformulation of Heegaard Floer homology. Geom. Topol. 10 (2006) 955-1096.

\bibitem{Yajing} Y. Liu. L-space surgeries on links. Quantum Topol. 8 (2017), no. 3, 505--570.

\bibitem{MM} P. Melvin, H. Morton. The coloured Jones function. Comm. Math. Phys. 169 (1995), no. 3, 501--520.



\bibitem{OSKnots} P. Ozsv\'ath, Z. Szab\'o. Holomorphic disks and knot invariants. Adv. Math. 186 (2004), no.1, 58--116. 



\bibitem{OSProperties} P. Ozsv\'ath, Z. Szab\'o. Holomorphic disks and three-manifold invariants: properties  and applications. Ann. of Math. (2) 159 (2004), no.3, 1159--1245. 

\bibitem{OSlens} P. Ozsv\'ath, Z. Szab\'o. On knot Floer homology and lens space surgeries. Topology 44 (2005), no.6, 1281--1300.


\bibitem{OSLinks} P. Ozsv\'ath, Z. Szab\'o. Holomorphic disks, link invariants and the multi-variable  Alexander polynomial. Algebr. Geom. Topol. 8 (2008), no. 2, 615--692.

\bibitem{RasmussenKnots} J. Rasmussen. Floer homology and knot complements. PhD thesis, Harvard University, 2003.



\bibitem{RasDiff} J. Rasmussen. Some differentials on Khovanov–Rozansky homology. Geometry \& Topology, 19(6), 3031--3104.

\bibitem{RT}  N. Y. Reshetikhin and V. G. Turaev. Ribbon graphs and their invariants derived from
quantum groups. Comm. Math. Phys. 127 (1990), 1--26.

\bibitem{RJ} M. Rosso and V. F. R. Jones. On the invariants of torus knots derived from quantum groups. Journal of Knot Theory and its 
Ramifications 2 (1993), 97--112.

\bibitem{Rozansky} L. Rozansky. An infinite torus braid yields a categorified Jones--Wenzl projector. Fundamenta Mathematicae 225 (2014): 305--326.

\bibitem{Wildi} A. Wildi. On Categorifications and Applications of Knot Invariants. PhD thesis, University of Z\"urich, 2025.

%\bibitem{Willis}

\bibitem{Zemke:graph} I. Zemke. Graph cobordisms and Heegaard Floer homology. arXiv:1512.01184.

\bibitem{Zemke} I. Zemke. Link cobordisms and functoriality in link Floer homology. J. Topol. 12 (2019), no. 1, 94--220. 

\bibitem{Zemke2} I. Zemke. Link cobordisms and absolute gradings on link Floer homology. Quantum Topol. 10 (2019), no. 2, 207--323. 

\bibitem{Zemke:basepoint} I. Zemke. Quasistabilization and basepoint moving maps in link Floer homology. Algebraic \& Geometric Topology 17 (2017), no. 6, 3461 -- 3518.

\bibitem{Zemke:bordered} I. Zemke. Bordered manifolds with torus boundary and the link surgery formula. arXiv:2109.11520.


\end{thebibliography}
\end{document}
